\documentclass[a4paper]{article}

\usepackage{graphicx}
\usepackage{amsmath,amssymb}
\usepackage{minted}
\usepackage{multirow}
\usepackage{float}
\usepackage{algorithm}
\usepackage{algpseudocode}
\usepackage{todonotes}
\usepackage{forest}
\usepackage{xstring}
\usepackage{xcolor}
\usepackage[most]{tcolorbox}
\usepackage{xparse}
\usepackage{natbib}

\usetikzlibrary{graphs,graphdrawing}
\usegdlibrary{layered}

% for external graphviz
\input{/home/donkarlo/Dropbox/repo/nd_language_ai_project/src/nd_language_ai/formal/computer/programming/tex/latex/code_clips/external_graphviz.tex}

% for tags
\input{/home/donkarlo/Dropbox/repo/nd_language_ai_project/src/nd_language_ai/formal/computer/programming/tex/latex/parts/control_sequence/kind/semtag/semtag.tex}

% for links
\input{/home/donkarlo/Dropbox/repo/nd_language_ai_project/src/nd_language_ai/formal/computer/programming/tex/latex/code_clips/seehere.tex}

\title{Novelty detection}
\date{}
\author{Mohammad Rahmani}

\begin{document}
   \maketitle

   Novelty detection is described as a neural function that determines whether an incoming stimulus is familiar or novel.
   The paper argues that such mechanisms are widespread in the central nervous system and participate in several cognitive functions. It proposes and simulates several neural novelty-detector models. A bi-compartmental model reproduces neural novelty-detection behavior particularly well and resembles cortical architecture, especially area 17. The model assumes that novelty detection begins with retrieving stored information and comparing it with incoming input \cite{salu-1988-models-of-neural-novelty-detectors-with-similarities-to-cerebral-cortex}

   Novelty detection is the mechanism by which an intelligent organism is able to identify an incoming sensory pattern as being so far unknown. If the pattern is sufficiently salient or associated with a high positive or strong negative utility, it will be given computational resources for effective future processing \seeherefooter{https://en.wikipedia.org/wiki/Novelty_detection}.

   \bibliographystyle{plainnat}
   \bibliography{/home/donkarlo/Dropbox/repo/nd_shared_research/bibliography/bibliography.bib}
\end{document}
